Over recent decades, both wired and wireless networks have undergone remarkable advancements, revolutionizing connectivity and enabling a wide array of use cases across industries and daily life. Wired networks have transitioned from copper-based infrastructure to advanced optical fiber technology, which offers significant advantages. Optical fibers, with their ability to transmit infrared light over long distances with minimal attenuation, provide gigabit-level speeds and serve as the backbone for critical infrastructures, including data centers and enterprise applications. Concurrently, wireless networks have experienced significant growth, with mobile networks evolving from 2G to the high-speed, low-latency capabilities of 5G, and Wi-Fi standards advancing from 802.11b to Wi-Fi 6 and the upcoming Wi-Fi 7 technologies.

In France, substantial efforts have been made to expand fiber and 5G networks. According to the French telecommunications regulator, \ac{ARCEP} \cite{arcepMarchxE9Communications}, as of the second quarter of 2024, 92\% of residents (32 million people) have access to high-speed wired connections, with 70\% (23 million) connected via \ac{FTTH}. Additionally, 5G adoption is rapidly increasing, with 18.5 million active SIM cards—representing 22\% of all mobile subscriptions—connected to 5G networks.

These advancements collectively enable networks to deliver impressive performance characterized by gigabit speeds, low latency, and high reliability. This progress is foundational for the emergence of innovative technologies and applications, including real-time remote collaboration, high-quality video streaming and gaming, immersive \ac{XR} experiences, autonomous systems, and the development of smart cities.

Among the many applications benefiting from advancements in network performance and the widespread deployment of multi-tier clouds, a category known as \ac{LL} applications has emerged. This class of applications, which includes \ac{CG}, cloud-based \ac{VR}, tele-robotics, and remote surgery, relies heavily on high-speed, \ac{LL} connections to deliver seamless user experiences.
Cloud gaming has revolutionized the gaming industry by shifting the computational burden from end-user devices to powerful cloud servers. This approach streams high-definition gameplay to users' devices, allowing even low-powered devices like smartphones and tablets to run graphically intensive games. \ac{CG} eliminates the need for expensive hardware and offers users access to a diverse library of games through subscription-based models, making gaming more accessible and affordable.

Although early \ac{CG} platforms, such as OnLive (launched in 2009), struggled due to limited network capacity, sparse cloud infrastructure, and insufficient device capabilities, recent advancements have changed the landscape. Today, industry giants like Google, Sony, Microsoft, NVIDIA, and Amazon have embraced \ac{CG}. These platforms leverage the network performance improvements and the global cloud infrastructure, positioning \ac{CG} as a fast-growing market. The global \ac{CG} market is projected to grow from USD 9.71 billion in 2024 to an astounding USD 126.62 billion by 2032 \cite{fortunebusinessinsightsCloudGaming}, reflecting its widespread adoption and potential.

Similarly, cloud-based \ac{VR} is driving the \ac{VR} market, which was valued at USD 28.78 billion in 2023 and is forecasted to reach USD 192.99 billion by 2032 \cite{nextmscVirtualReality}. By offloading intensive processing to the cloud, cloud \ac{VR} makes immersive experiences more accessible to both consumers and enterprises, eliminating the need for high-cost hardware. This democratization of \ac{VR} technology has paved the way for innovative use cases, including interactive virtual classrooms, enabling students to learn in immersive environments; medical training to provide lifelike simulations for surgeons to practice complex procedures; or military training, offering realistic simulations for combat and strategic planning, which reduce costs and risks compared to traditional methods.
Key players in the cloud \ac{VR} market include Meta Platforms, which has heavily invested in its vision of metaverses—shared virtual worlds where people interact as 3D avatars. Through its acquisition of Oculus \ac{VR}, Meta has developed a range of \ac{VR} hardware and software to support these initiatives. Other significant contributors to the cloud \ac{VR} ecosystem include NVIDIA, Samsung, Unity, Sony, and Microsoft, which continue to drive innovation and adoption across various sectors.

Despite the remarkable growth in market size and technological advancements, \ac{LL} applications still face significant challenges, particularly those arising from network variability. Mobile and Wi-Fi networks are prone to bandwidth fluctuations, increased delays, and jitter, all of which can jeopardize the user experience for latency-sensitive applications. For instance, in \ac{CG}, video quality and smooth gameplay depend on maintaining a stable downlink bitrate, low network delays, and minimal packet loss. Higher bitrates enable higher resolutions and frame rates, but they can also exacerbate delays due to phenomena like \textit{bufferbloat}, where excessive buffering in network queues leads to increased latency \cite{ronteix2022reduction}.

Moreover, current network architectures remain largely bandwidth-oriented, often neglecting the latency requirements of emerging applications. Latency challenges frequently arise from diverse sources, including network congestion, inefficient routing, and variable capacity links, making it difficult for applications to consistently deliver the expected \ac{QoE}. While advances like latency-based congestion control algorithms, optimized application-layer protocols, and innovative techniques such as \textit{negative latency} \cite{engadgetGoogleWants} have been proposed, their effectiveness is often limited by the underlying variability of time-varying capacity networks.

Although, solutions such as traffic prioritization mechanisms, including network slicing in 5G, offer promising avenues to address latency issues, these approaches raise concerns about net neutrality and the potential degradation of performance for non-prioritized traffic. These aforementioned challenges highlight the need for anomaly and diagnostic solutions for \ac{LL} applications.

%These challenges highlight the need for network operators to develop  novel solutions to detect and diagnose the possible causes of low-latency service degradations to provide network solutions and resources require for these applications, troubleshoot the performance issues to provide the best quality to clients or industrials.